\documentclass[a4paper,11pt]{report}

% =====================================================
% Packages
% =====================================================
\usepackage[utf8]{inputenc}
\usepackage[T1]{fontenc}
\usepackage[french]{babel}
\usepackage{geometry}
\usepackage{listings}
\usepackage{xcolor}
\usepackage{hyperref}

\geometry{margin=2.5cm}

% =====================================================
% Configuration listings C & Bash
% =====================================================
\definecolor{codegray}{rgb}{0.95,0.95,0.95}
\definecolor{darkgreen}{rgb}{0.0,0.5,0.0}
\definecolor{darkred}{rgb}{0.6,0.0,0.0}

\lstset{
  language=C,
  backgroundcolor=\color{codegray},
  basicstyle=\ttfamily\small,
  keywordstyle=\color{blue},
  commentstyle=\color{darkgreen},
  stringstyle=\color{darkred},
  frame=single,
  breaklines=true,
  tabsize=2,
  showstringspaces=false
}

% =====================================================
% Titre
% =====================================================
\title{\textbf{Qualité de code en C}\\
Bonnes pratiques, robustesse et compilation stricte}
\author{Document pédagogique complet}

\begin{document}

\maketitle
\tableofcontents
\newpage

% =====================================================
\chapter{Introduction}

\section{Objectif du document}

Ce document a pour objectif de présenter les bonnes pratiques essentielles
en langage C à travers une approche comparative entre \textbf{mauvais} et
\textbf{bon code}. Il vise à améliorer la qualité, la robustesse et la
maintenabilité des programmes en mettant en évidence les erreurs courantes
et leurs corrections.

\section{Criticité}

Le langage C est puissant mais peu permissif : il laisse au développeur la
responsabilité de la gestion mémoire, des types et du contrôle des erreurs.
Une mauvaise pratique peut rapidement conduire à :
\begin{itemize}
  \item des comportements indéfinis
  \item des failles de sécurité
  \item des crashs
  \item des bugs difficiles à détecter
\end{itemize}

Adopter des règles strictes permet de prévenir ces risques et de produire
un code fiable et industriel.

\section{Contenu du document}

Ce document couvre les aspects fondamentaux de la qualité en C :
\begin{itemize}
  \item Compilation stricte et gestion des warnings
  \item Bon usage des types et conversions
  \item Validation des entrées utilisateur
  \item Gestion des erreurs et des retours de fonctions
  \item Manipulation sûre de la mémoire (pointeurs, tableaux, allocation)
  \item Écriture de code lisible et maintenable (noms, commentaires, structure)
\end{itemize}

\section{Approche pédagogique}

Chaque chapitre suit une structure simple et efficace :
\begin{itemize}
  \item un exemple de \textbf{mauvais code}
  \item une version corrigée (\textbf{bon code})
  \item une \textbf{explication} claire du problème et de la solution
\end{itemize}

Cette approche permet de comprendre rapidement les erreurs fréquentes et
d'adopter les bons réflexes en situation réelle.

% =====================================================
\chapter{Compilation stricte (GCC)}

\section{Mauvais code}
\begin{lstlisting}
gcc fichier.c
\end{lstlisting}

\section{Bon code}
\begin{lstlisting}[language=bash]
gcc -std=c11 -Wall -Wextra -Werror \
    -Wshadow -Wconversion -Wsign-conversion \
    -Wstrict-prototypes -Wimplicit-function-declaration \
    fichier.c -o app
\end{lstlisting}

\section{Explication}
Une compilation stricte permet de détecter les erreurs très tôt. Avec \texttt{-Werror}, tous les warnings deviennent bloquants, ce qui garantit un code sans ambiguïté ni comportement indéfini.

L'ajout d'autres paramètres permet d'améliorer la commande gcc afin  d'y ajouter de nouvelles analyses. Cela permet ensuite au développeur de corriger et améliorer son programme plus efficacement.

% =====================================================
\chapter{Programme principal correct}

\section{Mauvais code}
\begin{lstlisting}
int main() {
    /* ... */
}
\end{lstlisting}

\section{Bon code}
\begin{lstlisting}
int main(void) {
    return 0;
}
\end{lstlisting}

\section{Explication}
En C, \texttt{main()} sans \texttt{void} signifie que la liste des arguments n'est pas spécifiée.
La forme \texttt{int main(void)} est plus claire, plus correcte et cohérente avec une compilation stricte.

% =====================================================
\chapter{Variables globales inutiles}

\section{Mauvais code}
\begin{lstlisting}
#include <stdio.h>

int count;

int main(void) {
    count = 5;
    printf("%d\n", count);
    return 0;
}
\end{lstlisting}

\section{Bon code}
\begin{lstlisting}
#include <stdio.h>

int main(void) {
    int count = 5;
    printf("%d\n", count);
    return 0;
}
\end{lstlisting}

\section{Explication}
Les variables globales augmentent le couplage entre fonctions et rendent le code plus difficile à tester et à maintenir.
Il faut préférer les variables locales tant qu'un partage global n'est pas justifié.

% =====================================================
\chapter{Variables non initialisées}

\section{Mauvais code}
\begin{lstlisting}
#include <stdio.h>

int main(void) {
    int x;
    if (x > 0) {
        printf("ok\n");
    }
    return 0;
}
\end{lstlisting}

\section{Bon code}
\begin{lstlisting}
#include <stdio.h>

int main(void) {
    int x = 0;
    if (x > 0) {
        printf("ok\n");
    }
    return 0;
}
\end{lstlisting}

\section{Explication}
Lire une variable non initialisée provoque un comportement indéfini.
Le compilateur peut parfois détecter ce cas, mais il faut surtout adopter une discipline systématique d'initialisation.

% =====================================================
\chapter{Valeurs de retour ignorées}

\section{Mauvais code}
\begin{lstlisting}
#include <stdio.h>

int main(void) {
    char buf[16];
    fgets(buf, sizeof buf, stdin);
    puts(buf);
    return 0;
}
\end{lstlisting}

\section{Bon code}
\begin{lstlisting}
#include <stdio.h>

int main(void) {
    char buf[16];

    if (fgets(buf, sizeof buf, stdin) == NULL) {
        return 1;
    }

    puts(buf);
    return 0;
}
\end{lstlisting}

\section{Explication}
Ignorer la valeur de retour d'une fonction d'entrée/sortie empêche de détecter un échec.
Il faut vérifier systématiquement les appels qui peuvent échouer.

% =====================================================
\chapter{Entrées utilisateur avec scanf}

\section{Mauvais code}
\begin{lstlisting}
#include <stdio.h>

int main(void) {
    int x;
    scanf("%d", &x);
    printf("%d\n", x);
    return 0;
}
\end{lstlisting}

\section{Bon code}
\begin{lstlisting}
#include <stdio.h>

int main(void) {
    int x = 0;

    if (scanf("%d", &x) != 1) {
        return 1;
    }

    printf("%d\n", x);
    return 0;
}
\end{lstlisting}

\section{Explication}
\texttt{scanf} peut échouer si l'entrée ne correspond pas au format attendu.
Sans test du retour, la variable peut rester indéterminée.

% =====================================================
\chapter{Conversion robuste avec strtoul}

\section{Mauvais code}
\begin{lstlisting}
#include <stdlib.h>

int main(void) {
    int x = atoi("12abc");
    return x;
}
\end{lstlisting}

\section{Bon code}
\begin{lstlisting}
#include <errno.h>
#include <stdlib.h>

int parse_u32(const char *s, unsigned long *out) {
    char *end = NULL;
    unsigned long v;

    errno = 0;
    v = strtoul(s, &end, 10);

    if (errno != 0 || end == s || *end != '\0') {
        return 1;
    }

    *out = v;
    return 0;
}
\end{lstlisting}

\section{Explication}
\texttt{atoi} ne permet pas de distinguer une vraie conversion d'une conversion partielle ou invalide.
\texttt{strtoul} associé à \texttt{errno} et au pointeur de fin permet une validation robuste.

% =====================================================
\chapter{Overflow signé}

\section{Mauvais code}
\begin{lstlisting}
#include <limits.h>

int add_bad(int a, int b) {
    return a + b;
}
\end{lstlisting}

\section{Bon code}
\begin{lstlisting}
#include <limits.h>

int add_good(int a, int b, int *out) {
    if (out == NULL) {
        return 1;
    }

    if ((b > 0 && a > INT_MAX - b) ||
        (b < 0 && a < INT_MIN - b)) {
        return 1;
    }

    *out = a + b;
    return 0;
}
\end{lstlisting}

\section{Explication}
Le dépassement de capacité d'un entier signé provoque un comportement indéfini.
Il faut vérifier les bornes avant l'opération.

% =====================================================
\chapter{Conversions signé / non signé}

\section{Mauvais code}
\begin{lstlisting}
#include <string.h>

int main(void) {
    int i = -1;
    size_t n = strlen("abc");
    return (i < n);
}
\end{lstlisting}

\section{Bon code}
\begin{lstlisting}
#include <string.h>

int main(void) {
    size_t i = 0;
    size_t n = strlen("abc");
    return (i < n);
}
\end{lstlisting}

\section{Explication}
Comparer un entier signé négatif avec un type non signé force une conversion implicite piégeuse.
Il faut choisir des types cohérents avec le domaine de valeurs attendu.

% =====================================================
\chapter{Masquage de variable}

\section{Mauvais code}
\begin{lstlisting}
int main(void) {
    int x = 10;

    if (x > 0) {
        int x = 3;
        (void)x;
    }

    return x;
}
\end{lstlisting}

\section{Bon code}
\begin{lstlisting}
int main(void) {
    int x = 10;

    if (x > 0) {
        int y = 3;
        (void)y;
    }

    return x;
}
\end{lstlisting}

\section{Explication}
Le shadowing rend le code confus et favorise les erreurs de lecture.
Il vaut mieux éviter de redéclarer une variable avec le même nom dans une portée interne.

% =====================================================
\chapter{Prototypes de fonctions}

\section{Mauvais code}
\begin{lstlisting}
int main(void) {
    return add(2, 3);
}

int add(int a, int b) {
    return a + b;
}
\end{lstlisting}

\section{Bon code}
\begin{lstlisting}
int add(int a, int b);

int main(void) {
    return add(2, 3);
}

int add(int a, int b) {
    return a + b;
}
\end{lstlisting}

\section{Explication}
Appeler une fonction sans prototype peut conduire à une signature supposée incorrecte.
En compilation stricte, chaque fonction doit être déclarée avant usage.

% =====================================================
\chapter{Strict prototypes}

\section{Mauvais code}
\begin{lstlisting}
int foo();
int main() {
    return foo();
}
\end{lstlisting}

\section{Bon code}
\begin{lstlisting}
int foo(void);

int main(void) {
    return foo();
}
\end{lstlisting}

\section{Explication}
En C, \texttt{int foo();} signifie que les arguments sont inconnus.
\texttt{int foo(void);} exprime correctement qu'il n'y a aucun argument.

% =====================================================
\chapter{Affectation au lieu de comparaison}

\section{Mauvais code}
\begin{lstlisting}
int main(void) {
    int x = 0;

    if (x = 1) {
        return 1;
    }

    return 0;
}
\end{lstlisting}

\section{Bon code}
\begin{lstlisting}
int main(void) {
    int x = 0;

    if (x == 1) {
        return 1;
    }

    return 0;
}
\end{lstlisting}

\section{Explication}
Utiliser \texttt{=} dans une condition modifie la variable au lieu de la comparer.
C'est une erreur classique que les warnings de compilation permettent souvent de repérer.

% =====================================================
\chapter{Switch incomplet}

\section{Mauvais code}
\begin{lstlisting}
void foo(void);

void f(int s) {
    switch (s) {
        case 1:
            foo();
    }
}
\end{lstlisting}

\section{Bon code}
\begin{lstlisting}
void foo(void);

void f(int s) {
    switch (s) {
        case 1:
            foo();
            break;
        default:
            break;
    }
}
\end{lstlisting}

\section{Explication}
Sans \texttt{break}, un cas peut tomber dans le suivant.
Sans \texttt{default}, un état non prévu n'est pas explicitement géré.

% =====================================================
\chapter{Accès hors bornes dans les tableaux}

\section{Mauvais code}
\begin{lstlisting}
int get_bad(const int *a, unsigned len, unsigned idx) {
    return a[idx];
}
\end{lstlisting}

\section{Bon code}
\begin{lstlisting}
int get_good(const int *a, unsigned len, unsigned idx, int *out) {
    if (a == NULL || out == NULL) {
        return 1;
    }

    if (idx >= len) {
        return 1;
    }

    *out = a[idx];
    return 0;
}
\end{lstlisting}

\section{Explication}
Un accès hors bornes provoque un comportement indéfini.
Il faut vérifier les bornes avant tout accès mémoire indexé.

% =====================================================
\chapter{Pointeurs nuls}

\section{Mauvais code}
\begin{lstlisting}
void set_bad(int *p) {
    *p = 42;
}
\end{lstlisting}

\section{Bon code}
\begin{lstlisting}
int set_good(int *p) {
    if (p == NULL) {
        return 1;
    }

    *p = 42;
    return 0;
}
\end{lstlisting}

\section{Explication}
Déréférencer un pointeur nul provoque généralement un crash.
Toute fonction manipulant un pointeur reçu de l'extérieur doit valider ce pointeur.

% =====================================================
\chapter{Allocation dynamique}

\section{Mauvais code}
\begin{lstlisting}
#include <stdlib.h>

int main(void) {
    int *p = malloc(10 * sizeof(int));
    p[0] = 1;
    free(p);
    return 0;
}
\end{lstlisting}

\section{Bon code}
\begin{lstlisting}
#include <stdlib.h>

int main(void) {
    int *p = malloc(10 * sizeof(*p));

    if (p == NULL) {
        return 1;
    }

    p[0] = 1;
    free(p);
    p = NULL;
    return 0;
}
\end{lstlisting}

\section{Explication}
\texttt{malloc} peut échouer et retourner \texttt{NULL}.
Après un \texttt{free}, remettre le pointeur à \texttt{NULL} réduit le risque d'usage accidentel ultérieur.

% =====================================================
\chapter{Erreur classique avec sizeof(pointer)}

\section{Mauvais code}
\begin{lstlisting}
#include <string.h>

int main(void) {
    const char *s = "hello";
    char buf[16];

    memcpy(buf, s, sizeof(s));
    return 0;
}
\end{lstlisting}

\section{Bon code}
\begin{lstlisting}
#include <string.h>

int main(void) {
    const char *s = "hello";
    char buf[16];

    memcpy(buf, s, strlen(s) + 1);
    return 0;
}
\end{lstlisting}

\section{Explication}
Quand \texttt{s} est un pointeur, \texttt{sizeof(s)} donne la taille du pointeur, pas celle de la chaîne pointée.
Il faut raisonner sur la donnée réelle, pas sur l'adresse.

% =====================================================
\chapter{Macros non sûres}

\section{Mauvais code}
\begin{lstlisting}
#define INC(x) x++
#define SWAP(a,b) { int t=a; a=b; b=t; }
\end{lstlisting}

\section{Bon code}
\begin{lstlisting}
#define INC(x) do { (x)++; } while (0)
#define SWAP(a,b) do { int t = (a); (a) = (b); (b) = t; } while (0)
\end{lstlisting}

\section{Explication}
Une macro sans parenthèses ni structure sûre peut casser la logique du code appelant.
Le motif \texttt{do \{ \} while (0)} est la forme classique pour une macro multi-instruction.

% =====================================================
\chapter{Usage inutile de goto}

\section{Mauvais code}
\begin{lstlisting}
int f(int ok) {
    if (!ok) {
        goto err;
    }

    return 0;

err:
    return 1;
}
\end{lstlisting}

\section{Bon code}
\begin{lstlisting}
int f(int ok) {
    if (!ok) {
        return 1;
    }

    return 0;
}
\end{lstlisting}

\section{Explication}
\texttt{goto} rend le flux d'exécution plus difficile à suivre.
Il peut avoir des usages légitimes dans certains schémas de nettoyage, mais il faut l'éviter quand une structure plus simple suffit.

% =====================================================
\chapter{Commentaires utiles}

\section{Mauvais code}
\begin{lstlisting}
i++; /* incremente i */
\end{lstlisting}

\section{Bon code}
\begin{lstlisting}
/* On limite les tentatives pour eviter une boucle infinie. */
retry_count++;
\end{lstlisting}

\section{Explication}
Un bon commentaire explique l'intention ou la contrainte, pas une évidence syntaxique.
Le lecteur doit comprendre pourquoi le code existe.

% =====================================================
\chapter{Paramètre inutilisé}

\section{Mauvais code}
\begin{lstlisting}
void cb(int x) {
}
\end{lstlisting}

\section{Bon code}
\begin{lstlisting}
void cb(int x) {
    (void)x;
}
\end{lstlisting}

\section{Explication}
Un paramètre non utilisé déclenche souvent un warning.
Le cast en \texttt{(void)} documente explicitement que ce paramètre est volontairement ignoré.

% =====================================================
\chapter{Nommage cohérent}

\section{Mauvais code}
\begin{lstlisting}
#define max_retry 3
struct Person { int Age; };
int DoStuff();
\end{lstlisting}

\section{Bon code}
\begin{lstlisting}
#define MAX_RETRY 3

typedef struct {
    int age;
} person_t;

int do_stuff(void);
\end{lstlisting}

\section{Explication}
Un nommage cohérent améliore fortement la lisibilité.
Une convention simple peut être :
\begin{itemize}
    \item macros en majuscules ;
    \item types structurés avec suffixe \texttt{\_t} si le projet l'autorise ;
    \item fonctions et variables en minuscules avec underscore.
\end{itemize}

% =====================================================
\chapter{Style de code cohérent}

\section{Mauvais code}
\begin{lstlisting}
if(x>0){foo();}
\end{lstlisting}

\section{Bon code}
\begin{lstlisting}
if (x > 0) {
    foo();
}
\end{lstlisting}

\section{Explication}
Un style homogène rend le code plus facile à lire, relire et maintenir.
Peu importe la convention choisie, l'essentiel est d'être cohérent dans tout le projet.

% =====================================================
\chapter{Conclusion}

Un code C de qualité professionnelle repose sur :
\begin{itemize}
    \item une compilation stricte sans warning ;
    \item une gestion rigoureuse des types, des conversions et des prototypes ;
    \item des vérifications systématiques des retours, pointeurs et bornes ;
    \item un style cohérent et un nommage lisible ;
    \item une attention particulière à la robustesse mémoire et à la sécurité.
\end{itemize}

Ces pratiques permettent de produire un code plus fiable, plus maintenable et plus sûr dans la durée.

\end{document}